% CVPR 2026 Paper Template; see https://github.com/cvpr-org/author-kit

\documentclass[10pt,twocolumn,letterpaper]{article}

%%%%%%%%% PAPER TYPE  - PLEASE UPDATE FOR FINAL VERSION
\usepackage{cvpr}              % To produce the CAMERA-READY version
% \usepackage[review]{cvpr}      % To produce the REVIEW version
% \usepackage[pagenumbers]{cvpr} % To force page numbers, e.g. for an arXiv version

% Import additional packages in the preamble file, before hyperref
\input{preamble}

% It is strongly recommended to use hyperref, especially for the review version.
% hyperref with option pagebackref eases the reviewers' job.
% Please disable hyperref *only* if you encounter grave issues, 
% e.g. with the file validation for the camera-ready version.
%
% If you comment hyperref and then uncomment it, you should delete *.aux before re-running LaTeX.
% (Or just hit 'q' on the first LaTeX run, let it finish, and you should be clear).
\definecolor{cvprblue}{rgb}{0.21,0.49,0.74}
\usepackage[pagebackref,breaklinks,colorlinks,allcolors=cvprblue]{hyperref}

%%%%%%%%% PAPER ID  - PLEASE UPDATE
% \def\paperID{*****} % *** Enter the Paper ID here
\def\confName{CVPR}
\def\confYear{2026}

%%%%%%%%% TITLE - PLEASE UPDATE
\title{GENERATIVE OPTICAL OBFUSCATION THROUGH NATURAL EXPOSURE RECONSTRUCTION}

%%%%%%%%% AUTHORS - PLEASE UPDATE
\author{Sodi Kroehler\\
University of Pittsburgh\\
{\tt\small sek188@pitt.edu}
% For a paper whose authors are all at the same institution,
% omit the following lines up until the closing ``}''.
% Additional authors and addresses can be added with ``\and'',
% just like the second author.
% To save space, use either the email address or home page, not both
\and
Tanmay Shetty\\
University of Pittsburgh\\
% {\tt\small sek188@pitt.edu}
}

\begin{document}
\maketitle


\begin{abstract}
The ABSTRACT is to be in fully justified italicized text, at the top of the left-hand column, below the author and affiliation information.
Use the word ``Abstract'' as the title, in 12-point Times, boldface type, centered relative to the column, initially capitalized.
The abstract is to be in 10-point, single-spaced type.
Leave two blank lines after the Abstract, then begin the main text.
Look at previous \confName abstracts to get a feel for style and length.
\end{abstract}


\input{writeup/parts/intro}
\input{writeup/parts/previous_work}
\input{writeup/parts/data}
\input{writeup/parts/meth}
\input{writeup/parts/results}
\input{writeup/parts/conclusion}
{
    \small
    \bibliographystyle{ieeenat_fullname}
    \bibliography{main}
}

% WARNING: do not forget to delete the supplementary pages from your submission 
% \input{sec/X_suppl}

\end{document}
